% GENERATED FILE. Edit canonical lessons, then run npm run generate:books.
% canonical-source-hash: fnv1a64:b03b977fdc1f7922
% canonical-lessons: KA-C241-tikisu, KA-C241-hancikollu, KA-C241-dvesisu, KA-C241-kanduhidi, KA-C241-sala-pade

\chapter{To criticise, To share, To hate, To find out, To borrow}
\label{ch:ka-to-criticise-to-share-to-hate-to-find-out-to-borrow}
\hlchaptermodality{241}

\begin{chapteropening}
\textbf{By the end of this chapter:} \emph{I can say to criticise, to share, to hate, to find out and to borrow.}

Complete the last lesson of chapter 241: I can say to criticise, to share, to hate, to find out and to borrow.
\end{chapteropening}

\section[\texorpdfstring{\kn{ಟೀಕಿಸು} (ṭīkisu)}{ಟೀಕಿಸು (ṭīkisu)}]{\kn{ಟೀಕಿಸು} (ṭīkisu) — to criticise}

\label{lesson:KA-C241-tikisu}



\begin{quote}
Before the new one: say the Kannada for to argue, then the Kannada for to imagine.
\end{quote}

\subsection*{You'll want to know: \kn{ಟೀಕಿಸು}}
\textbf{\kn{ಟೀಕಿಸು}} — \emph{ṭīkisu} — \textquotedblleft{}to criticise\textquotedblright{}.

\begin{grammarlens}[title={What you've built}]
One more word for this chapter.
\end{grammarlens}

\subsection*{Your turn}
\begin{itemize}
\raggedright
  \item \emph{Say it:} \emph{ṭīkisu}
  \item \emph{Say it:} \emph{ṭīkisu}, once more
  \item \emph{Say it:} say \emph{ṭīkisu}
  \item \emph{Recall:} say the Kannada for to argue, then the Kannada for to imagine, then say \emph{ṭīkisu} again
\end{itemize}

\begin{tcolorbox}[breakable,colback=teal!4,colframe=teal!35!black,title={Before you move on}]
What does \kn{ಟೀಕಿಸು} mean? (\textquotedblleft{}To criticise\textquotedblright{}.) Say it once more.
\end{tcolorbox}
\section[\texorpdfstring{\kn{ಹಂಚಿಕೊಳ್ಳು} (hañcikoḷḷu)}{ಹಂಚಿಕೊಳ್ಳು (hañcikoḷḷu)}]{\kn{ಹಂಚಿಕೊಳ್ಳು} (hañcikoḷḷu) — to share}

\label{lesson:KA-C241-hancikollu}



\begin{quote}
Before the new one: say the Kannada for to imagine, then the Kannada for to criticise.
\end{quote}

\subsection*{You'll want to know: \kn{ಹಂಚಿಕೊಳ್ಳು}}
\textbf{\kn{ಹಂಚಿಕೊಳ್ಳು}} — \emph{hañcikoḷḷu} — \textquotedblleft{}to share\textquotedblright{}.

\begin{grammarlens}[title={What you've built}]
One more word for this chapter.
\end{grammarlens}

\subsection*{Your turn}
\begin{itemize}
\raggedright
  \item \emph{Say it:} \emph{hañcikoḷḷu}
  \item \emph{Say it:} \emph{hañcikoḷḷu}, once more
  \item \emph{Say it:} say \emph{hañcikoḷḷu}
  \item \emph{Recall:} say the Kannada for to imagine, then the Kannada for to criticise, then say \emph{hañcikoḷḷu} again
\end{itemize}

\begin{tcolorbox}[breakable,colback=teal!4,colframe=teal!35!black,title={Before you move on}]
What does \kn{ಹಂಚಿಕೊಳ್ಳು} mean? (\textquotedblleft{}To share\textquotedblright{}.) Say it once more.
\end{tcolorbox}
\section[\texorpdfstring{\kn{ದ್ವೇಷಿಸು} (dvēṣisu)}{ದ್ವೇಷಿಸು (dvēṣisu)}]{\kn{ದ್ವೇಷಿಸು} (dvēṣisu) — to hate}

\label{lesson:KA-C241-dvesisu}



\begin{quote}
Before the new one: say the Kannada for to criticise, then the Kannada for to share.
\end{quote}

\subsection*{You'll want to know: \kn{ದ್ವೇಷಿಸು}}
\textbf{\kn{ದ್ವೇಷಿಸು}} — \emph{dvēṣisu} — \textquotedblleft{}to hate\textquotedblright{}.

\begin{grammarlens}[title={What you've built}]
One more word for this chapter.
\end{grammarlens}

\subsection*{Your turn}
\begin{itemize}
\raggedright
  \item \emph{Say it:} \emph{dvēṣisu}
  \item \emph{Say it:} \emph{dvēṣisu}, once more
  \item \emph{Say it:} say \emph{dvēṣisu}
  \item \emph{Recall:} say the Kannada for to criticise, then the Kannada for to share, then say \emph{dvēṣisu} again
\end{itemize}

\begin{tcolorbox}[breakable,colback=teal!4,colframe=teal!35!black,title={Before you move on}]
What does \kn{ದ್ವೇಷಿಸು} mean? (\textquotedblleft{}To hate\textquotedblright{}.) Say it once more.
\end{tcolorbox}
\section[\texorpdfstring{\kn{ಕಂಡುಹಿಡಿ} (kaṇḍuhiḍi)}{ಕಂಡುಹಿಡಿ (kaṇḍuhiḍi)}]{\kn{ಕಂಡುಹಿಡಿ} (kaṇḍuhiḍi) — to find out, to discover}

\label{lesson:KA-C241-kanduhidi}



\begin{quote}
Before the new one: say the Kannada for to share, then the Kannada for to hate.
\end{quote}

\subsection*{You'll want to know: \kn{ಕಂಡುಹಿಡಿ}}
\textbf{\kn{ಕಂಡುಹಿಡಿ}} — \emph{kaṇḍuhiḍi} — \textquotedblleft{}to find out, to discover\textquotedblright{}.

\begin{grammarlens}[title={What you've built}]
One more word for this chapter.
\end{grammarlens}

\subsection*{Your turn}
\begin{itemize}
\raggedright
  \item \emph{Say it:} \emph{kaṇḍuhiḍi}
  \item \emph{Say it:} \emph{kaṇḍuhiḍi}, once more
  \item \emph{Say it:} say \emph{kaṇḍuhiḍi}
  \item \emph{Recall:} say the Kannada for to share, then the Kannada for to hate, then say \emph{kaṇḍuhiḍi} again
\end{itemize}

\begin{tcolorbox}[breakable,colback=teal!4,colframe=teal!35!black,title={Before you move on}]
What does \kn{ಕಂಡುಹಿಡಿ} mean? (\textquotedblleft{}To find out, to discover\textquotedblright{}.) Say it once more.
\end{tcolorbox}
\section[\texorpdfstring{\kn{ಸಾಲ} \kn{ಪಡೆ} (sāla paḍe)}{ಸಾಲ ಪಡೆ (sāla paḍe)}]{\kn{ಸಾಲ} \kn{ಪಡೆ} (sāla paḍe) — to borrow (money)}

\label{lesson:KA-C241-sala-pade}



\begin{quote}
Before the new one: say the Kannada for to hate, then the Kannada for to find out.
\end{quote}

\subsection*{You'll want to know: \kn{ಸಾಲ} \kn{ಪಡೆ}}
\textbf{\kn{ಸಾಲ} \kn{ಪಡೆ}} — \emph{sāla paḍe} — \textquotedblleft{}to borrow (money)\textquotedblright{}.

\begin{grammarlens}[title={What you've built}]
One more word for this chapter.
\end{grammarlens}

\subsection*{Your turn}
\begin{itemize}
\raggedright
  \item \emph{Say it:} \emph{sāla paḍe}
  \item \emph{Say it:} \emph{sāla paḍe}, once more
  \item \emph{Say it:} say \emph{sāla paḍe}
  \item \emph{Recall:} say the Kannada for to hate, then the Kannada for to find out, then say \emph{sāla paḍe} again
\end{itemize}

\begin{tcolorbox}[breakable,colback=teal!4,colframe=teal!35!black,title={Before you move on}]
What does \kn{ಸಾಲ} \kn{ಪಡೆ} mean? (\textquotedblleft{}To borrow (money)\textquotedblright{}.) Say it once more.
\end{tcolorbox}
